\chapter{The Second Variation}
% OUTLINE Ch2 "The Second Variation (topology's face: the obstruction
% structure)". Laws: INTUITION LAW (labeled blocks, adds-no-claim),
% CITATION (Legendre 1786 condition; Jacobi 1837 conjugate points; Morse
% 1934 The Calculus of Variations in the Large), LENS EXHIBIT 1 OF 2
% rides here under the standing two fences (verb = illustrates; the
% obstruction/topology NAME earned-not-asserted; "the SAME residue as
% Vol 2 Ch6's exhibit, now worn as obstruction-shape, same calibration"
% -- Beastmaster's words, Beastmaster reads on arrival). Gate: Kodo.

\begin{intuition}
Stand at a mountain pass. Underfoot it is as flat as any valley floor ---
a marble placed exactly there would rest --- and yet the flatness lies.
Step along the ridge and you fall away to either side; step across it and
you climb. The valley floor and the pass agree completely at first order
and differ in everything that matters, and the difference is invisible to
any instrument that only asks whether the ground is level. The feeling to
carry: flatness at first order hides the real story, and the real story
is told by what happens at second order --- and, deeper, by the shape of
the country that made a pass necessary at all.
\end{intuition}

The first half of the story is classical and carries two names. The last
chapter's structure finds the flat places --- first variation zero --- and
stops; to ask which flat places are minima, the probe must go one order
deeper. Vary the candidate twice: the second variation is a quadratic
reading of the wiggle, and its sign pattern classifies the flat place ---
positive against every wiggle, a floor; negative against every wiggle, a
crest; mixed, a pass. Legendre, in 1786, gave the pointwise condition the
sign demands of the integrand, and Jacobi, in 1837, found the deeper and
stranger fact: follow a family of extremals out from a point, and there
can come a definite station --- the conjugate point --- past which a path
that minimized honestly stops minimizing, its second variation acquiring
a direction of descent. The reader has met the physical shadow of this on
a globe: the great-circle arc from a city is genuinely shortest --- until
it passes the antipode, its conjugate point, after which a shorter way
around exists. The second variation is not a bookkeeping refinement; it
is where the geometry of the \emph{space of candidates} first speaks.

And that is the door to this chapter's real face. Ask the naive question:
could the pass be gotten rid of? Flatten it, push it down, deform the
landscape until only floors remain? For an isolated function on a
featureless space, perhaps. But on a country with \emph{shape} --- and
here the reader should picture walking on the surface of a doughnut ---
the answer is no, and the no is a theorem. Assign heights to the
doughnut's surface any smooth way you please: there must be a lowest
point and a highest, and there must \emph{also} be passes --- at least
two of them --- because the surface's two independent ways of looping
cannot be filled in by any rearrangement of heights. Marston Morse, in
the work gathered in his 1934 \emph{The Calculus of Variations in the
Large}, made this exact bargain precise: the critical points of a
function, each counted with its index --- the number of independent
downhill directions --- are constrained from below by the topology of
the space they live on. The holes demand the passes. Deformation relocates a
pass, reshapes it, slides it along the ridge; what deformation cannot do
is delete it, because the demand does not come from the function --- it
comes from the space.

Say the proof structure plainly, because it is unlike the last chapter's
and the contrast is the lesson. The extremal structure proves a
candidate best by varying it. The obstruction structure proves a thing
\emph{exists} without ever constructing it: read the topology, count
what it demands, and conclude that the landscape must contain a pass ---
somewhere, in some form, beyond the reach of any deformation to remove.
It is proof by inevitability. The quantity the argument rests on is an
invariant --- a number no continuous rearrangement can change --- and
the conclusion is not ``here is the pass'' but ``there is no world
without one.'' Where the first variation's signature was one arbitrary
wiggle carrying a universal quantifier, the second face's signature is
an invariant carrying an existential one.

\begin{intuition}
The feeling for it: a bump under a carpet. Press it flat here and it
rises there; chase it around the room all day and the one thing that
never changes is that there is a bump. The carpet's excess is not a
property of any place --- it is a property of the carpet --- and the
chasing is how you discover that fact, not how you escape it.
\end{intuition}

\begin{fence}
An illustration, not an identity --- and this volume's first of two
device appearances. The reader of this series has already met, in the
second volume's curvature chapter, the apparatus whose self-describing
loop returns changed: a residue that survives, reproducibly, however the
descent is refined. That residue is offered here, again and under the
same rules, as an illustration of \emph{shape} --- this chapter's shape:
a quantity that deformation relocates but cannot delete, the way a
demanded pass slides along a ridge but never leaves the landscape. It is
the same residue, worn now as obstruction-shape rather than loop-shape,
and the same calibration binds: the device does not thereby do Morse
theory, no invariant of any space has been computed by it, and whether
its residue could ever be \emph{earned} as an obstruction in the cited
sense stays on the far side of the series' standing fence. Illustration,
not identity.
\end{fence}

What the obstruction face cannot say is how much the surviving thing
weighs. Topology demands the pass but not its height; the invariant
forces existence and is silent about magnitude. When a quantity is not
merely forced to exist but forced to \emph{cancel exactly} --- when two
tellings of one till agree to the last coin, and the zero is a theorem
rather than a coincidence --- the proof at work has a different
structure again, and it is the oldest one in the ledger books. That is
algebra's face, and the next chapter turns to it.
